\documentclass[11pt,a4paper]{article}

\usepackage[utf8]{inputenc}
\usepackage[T1]{fontenc}
\usepackage{amsmath,amssymb,amsthm}
\usepackage{geometry}
\usepackage{booktabs}
\usepackage{hyperref}
\usepackage{xcolor}

\geometry{margin=2.5cm}

\hypersetup{
    colorlinks=true,
    linkcolor=blue!60!black,
    urlcolor=blue!60!black
}

\newtheorem{definition}{Definição}
\newtheorem{remark}{Observação}

\title{Laplaciano em Grafos aplicado ao pipeline DINOv2 + KNN}
\author{Trabalho de Teoria Espectral}
\date{\today}

\begin{document}

\maketitle

\section{Objetivo}

Este documento descreve, passo a passo, a matemática implementada no código
\path{Code/Core/Laplaciano_Propagacao/propagando_rotulos.py}.
O pipeline parte de imagens, gera embeddings com DINOv2, constrói um grafo KNN
ponderado e aplica um método espectral baseado no Laplaciano normalizado.

No código atual, a configuração fixa principal é:
\[
\texttt{DATASET\_NAME}=\texttt{CUB\_Cleaned50}, \quad
\texttt{K}=20, \quad
\texttt{NUM\_CLUSTERS}=50.
\]

\begin{remark}
Apesar do nome ``propagação de rótulos'', o código atual realiza clustering
espectral não supervisionado. Os rótulos reais entram apenas no final, para
avaliar a qualidade dos clusters encontrados.
\end{remark}

\section{Dados e embeddings}

Considere um conjunto de imagens
\[
\mathcal{I}=\{I_1,I_2,\ldots,I_n\}.
\]
No caso atual do CUB limpo,
\[
n=2889.
\]

Cada imagem é passada por uma rede DINOv2. O resultado é um vetor de atributos
\[
x_i \in \mathbb{R}^{384}.
\]
Agrupando todos os vetores, obtemos a matriz de embeddings
\[
X =
\begin{bmatrix}
--- x_1^\top ---\\
--- x_2^\top ---\\
\vdots\\
--- x_n^\top ---
\end{bmatrix}
\in \mathbb{R}^{n\times 384}.
\]

No código, esta matriz é salva em:
\[
\texttt{DataSets/Emb/CUB\_Cleaned50\_emb.pt}
\]
e também em formato NumPy:
\[
\texttt{DataSets/Emb/CUB\_Cleaned50\_emb.npy}.
\]

\section{Listas ranqueadas por vizinhança}

Para cada imagem \(i\), o código usa uma estrutura \texttt{BallTree} para
encontrar os vizinhos mais próximos no espaço dos embeddings.

Assim, para cada \(i\), temos uma lista ranqueada:
\[
R_i = (i, r_{i,1}, r_{i,2}, \ldots, r_{i,m}),
\]
onde o primeiro elemento é normalmente a própria imagem, e
\(r_{i,1}\) é o vizinho mais próximo de \(i\), \(r_{i,2}\) é o segundo mais
próximo, e assim por diante.

Essas listas são salvas no arquivo:
\[
\texttt{DataSets/Runs/CUB\_Cleaned50\_dinov2\_vits14\_output.json}.
\]

\section{Construção do grafo KNN ponderado}

\begin{definition}[Grafo ponderado]
Um grafo ponderado é uma tripla
\[
G=(V,E,W),
\]
onde \(V\) é o conjunto de vértices, \(E\) é o conjunto de arestas e
\(W\) contém pesos não negativos associados às arestas.
\end{definition}

No nosso caso, cada imagem é um vértice:
\[
V=\{1,2,\ldots,n\}.
\]

O código escolhe os \(K\) primeiros vizinhos de cada lista ranqueada. Para
\(K=20\), cada imagem escolhe seus 20 vizinhos mais próximos.

Construímos primeiro uma matriz direcionada \(W^{dir}\in\mathbb{R}^{n\times n}\).
Se \(j\) aparece na posição \(r\) da lista de vizinhos de \(i\), então:
\[
W^{dir}_{ij} = \frac{1}{r}.
\]
Logo, o primeiro vizinho recebe peso \(1\), o segundo recebe peso \(1/2\), o
terceiro recebe peso \(1/3\), etc.

No código:
\[
\texttt{W[i,j] = max(W[i,j], 1.0/rank)}.
\]

Depois, o grafo é tornado não direcionado por simetrização:
\[
W = W^{dir} + (W^{dir})^\top.
\]

Essa escolha faz com que relações recíprocas fiquem mais fortes. Por exemplo,
se \(i\) escolhe \(j\) com peso \(1\), e \(j\) escolhe \(i\) com peso \(1/2\),
então:
\[
W_{ij}=W_{ji}=1+\frac{1}{2}=1.5.
\]

\section{Grau dos vértices}

O grau ponderado de um vértice \(i\) é a soma dos pesos das arestas ligadas a
ele:
\[
d_i = \sum_{j=1}^{n} W_{ij}.
\]

A matriz diagonal dos graus é:
\[
D =
\begin{bmatrix}
d_1 & 0 & \cdots & 0\\
0 & d_2 & \cdots & 0\\
\vdots & \vdots & \ddots & \vdots\\
0 & 0 & \cdots & d_n
\end{bmatrix}.
\]

No código, não é necessário montar \(D\) explicitamente. Basta calcular o vetor
\[
d^{-1/2}_i = \frac{1}{\sqrt{d_i}}.
\]

\section{Laplaciano normalizado}

O Laplaciano normalizado usado no código é:
\[
L = I - D^{-1/2}WD^{-1/2}.
\]

Em coordenadas, a matriz normalizada de adjacência é:
\[
S_{ij} = \frac{W_{ij}}{\sqrt{d_i d_j}}.
\]
Assim,
\[
L = I - S.
\]

No código, isto aparece como:
\[
\texttt{normalized\_adjacency = d\_inv\_sqrt[:, None] * W * d\_inv\_sqrt[None, :]}
\]
e depois:
\[
\texttt{L = np.eye(n\_samples) - normalized\_adjacency}.
\]

\begin{remark}
A normalização por \(D^{-1/2}\) reduz o efeito de vértices com grau muito alto.
Assim, o espectro do grafo passa a refletir melhor a estrutura de comunidades,
e não apenas a quantidade de conexões de cada vértice.
\end{remark}

\section{Autovetores e embedding espectral}

A ideia central é usar os autovetores associados aos menores autovalores do
Laplaciano. Calculamos:
\[
L u_\ell = \lambda_\ell u_\ell,
\]
com
\[
\lambda_1 \leq \lambda_2 \leq \cdots \leq \lambda_c,
\]
onde \(c=\texttt{NUM\_CLUSTERS}\).

No caso atual:
\[
c=50.
\]

Formamos a matriz
\[
U =
\begin{bmatrix}
| & | & & |\\
u_1 & u_2 & \cdots & u_c\\
| & | & & |
\end{bmatrix}
\in \mathbb{R}^{n\times c}.
\]

Cada linha de \(U\) representa uma imagem no espaço espectral:
\[
U_{i,:} \in \mathbb{R}^{c}.
\]

\section{Normalização das linhas}

Seguindo a forma usual do algoritmo de Ng--Jordan--Weiss, o código normaliza
cada linha de \(U\):
\[
Y_{i,:} = \frac{U_{i,:}}{\|U_{i,:}\|_2}.
\]

Assim, a matriz final usada pelo KMeans é:
\[
Y \in \mathbb{R}^{n\times c}.
\]

Intuitivamente, o KMeans passa a comparar mais a direção dos pontos no espaço
espectral do que o tamanho dos vetores.

\section{Clustering com KMeans}

O KMeans é aplicado nas linhas de \(Y\), não nos embeddings originais do DINOv2:
\[
\text{cluster}(i) = \operatorname{KMeans}(Y_{i,:}).
\]

O resultado é um vetor:
\[
z \in \{0,1,\ldots,c-1\}^{n},
\]
onde \(z_i\) é o cluster atribuído à imagem \(i\).

\section{Rótulos reais e avaliação}

Para o dataset \texttt{CUB\_Cleaned50}, os rótulos reais são lidos de:
\[
\texttt{DataSets/CUB\_Cleaned50/manifest.csv}.
\]

O CUB usa classes numeradas de 1 até 50. No código, elas são convertidas para:
\[
0,1,\ldots,49.
\]

Seja \(y_i\) o rótulo real da imagem \(i\). Como os nomes dos clusters do
KMeans são arbitrários, não podemos comparar diretamente \(z_i=y_i\).

Por isso, construímos uma matriz de confusão:
\[
C_{ab} = \#\{i : z_i=a \ \text{e} \ y_i=b\}.
\]

Depois usamos o algoritmo Húngaro para encontrar o melhor pareamento um-para-um:
\[
\pi(a)=b.
\]

A acurácia reportada é:
\[
\text{acc} =
\frac{1}{n}
\sum_{i=1}^{n}
\mathbf{1}\{\pi(z_i)=y_i\}.
\]

Além disso, o código calcula:
\[
\text{ARI} = \text{Adjusted Rand Index},
\]
e
\[
\text{NMI} = \text{Normalized Mutual Information}.
\]

\section{Arquivos gerados}

O script salva os principais resultados em \texttt{DataSets/Laplacian}.
Com a configuração atual, o prefixo de saída é:
\[
\texttt{CUB\_Cleaned50\_laplacian\_label\_propagation\_k20\_c50}.
\]

\begin{center}
\begin{tabular}{ll}
\toprule
Arquivo & Conteúdo\\
\midrule
\texttt{\_clusters.npy} & Cluster bruto do KMeans para cada imagem\\
\texttt{\_mapped\_labels.npy} & Cluster convertido para o melhor rótulo real\\
\texttt{\_embedding.npy} & Matriz espectral \(Y\)\\
\texttt{\_eigenvalues.npy} & Autovalores escolhidos do Laplaciano\\
\texttt{\_confusion.npy} & Matriz de confusão cluster \(\times\) classe\\
\texttt{.json} & Grafo com posições, arestas, clusters e acertos\\
\texttt{.png} & Grafo colorido por cluster\\
\texttt{\_correctness.png} & Grafo colorido por acerto/erro\\
\texttt{\_embedding.png} & Projeção das duas primeiras coordenadas espectrais\\
\bottomrule
\end{tabular}
\end{center}

\section{Resumo do algoritmo}

\begin{enumerate}
    \item Ler as listas ranqueadas de vizinhos do arquivo JSON.
    \item Construir a matriz ponderada \(W^{dir}\) usando pesos \(1/r\).
    \item Simetrizar a matriz: \(W = W^{dir} + (W^{dir})^\top\).
    \item Calcular os graus \(d_i=\sum_j W_{ij}\).
    \item Construir o Laplaciano normalizado:
    \[
    L=I-D^{-1/2}WD^{-1/2}.
    \]
    \item Calcular os \(c\) menores autovetores de \(L\).
    \item Montar a matriz espectral \(U\) e normalizar suas linhas, obtendo \(Y\).
    \item Rodar KMeans nas linhas de \(Y\).
    \item Ler os rótulos reais do manifest.
    \item Usar o algoritmo Húngaro para mapear clusters para classes.
    \item Salvar métricas, arrays, JSON e figuras.
\end{enumerate}

\section{Comentário final}

O código implementa um método espectral sobre um grafo de similaridade entre
imagens. O DINOv2 fornece uma representação inicial das imagens; o KNN converte
essa representação em um grafo; o Laplaciano extrai a estrutura global do grafo;
e o KMeans transforma as coordenadas espectrais em clusters discretos.

\end{document}
